\paragraph{Notation.}
For \(r\geq1\), write \([r]:=\{1,\ldots,r\}\).  We use bold lowercase
Roman or Greek letters for vectors and bold uppercase letters for matrices.
For a vector \(\mathbf v\), its \(k\)-th coordinate is denoted by \(v_k\).
For a matrix \(\mathbf A\), its entry in row \(i\) and column \(j\) is denoted
by \(A_{ij}\).  Vector inequalities are coordinatewise.  For
\(1\leq p\leq\infty\),
\(\|\mathbf v\|_p\) is the usual \(\ell_p\)-norm, while
\(\|\mathbf v\|_0:=|\{i:v_i\neq0\}|\) is the
\emph{support size}.  We write
\(\mathbf 0_r\) and \(\mathbf 1_r\) for the length-\(r\) all-zero and
all-ones vectors.  For scalars, finite sets,
and binary strings, \(|\cdot|\) denotes absolute value, cardinality, and bit
length, respectively.

\subsection{Lattices}
\label{sec:simple-lattices}

For integers \(m\geq n\geq1\), let
\(\mathbf B\in\mathbb{R}^{m\times n}\) have linearly independent columns. The
\emph{lattice generated by \(\mathbf B\)} is defined by
\[
  \mathcal{L}(\mathbf B)\mathrel{:=}\mathbf B\mathbb{Z}^n
  =\{\mathbf B\mathbf z:\mathbf z\in\mathbb{Z}^n\}.
\]
The columns of $\mathbf B$ form a \emph{basis} of
$\mathcal{L}(\mathbf B)$. If $m=n$, then $\mathcal{L}(\mathbf B)$ is called
\emph{full rank}.

For a lattice \(\mathcal L\subseteq\mathbb R^m\) and
\(\mathbf t\in\mathbb R^m\), define
\[
  \operatorname{dist}(\mathbf t,\mathcal{L})
  \mathrel{:=}\min_{\mathbf v\in\mathcal{L}}\|\mathbf t-\mathbf v\|_2.
\]

\begin{definition}[Exact CVP]
Given a rational lattice basis $\mathbf B$, a rational target $\mathbf t$,
and a nonnegative rational squared threshold $\Delta$, exact decision-CVP
asks whether
\[
  \dist\bigl(\mathbf t,\mathcal L(\mathbf B)\bigr)^2\leq\Delta.
\]
Given $\mathbf B$ and $\mathbf t$, exact search-CVP asks for a vector
$\mathbf v\in\mathcal L(\mathbf B)$ satisfying
\[
  \|\mathbf t-\mathbf v\|_2
  =\dist\bigl(\mathbf t,\mathcal L(\mathbf B)\bigr).
\]
\end{definition}
The exact shortest vector problem (SVP) is the corresponding problem with
$\mathbf t=\mathbf 0_m$, except that the zero lattice vector is excluded.

\ifdefined\withModuleSIVP
\begin{definition}[Successive minima]
Let \(\mathcal L\subseteq\mathbb R^m\) be a lattice of rank \(n\), and let
\(\mathcal B_2^m:=\{\mathbf x\in\mathbb R^m:\|\mathbf x\|_2\leq1\}\) be the
closed Euclidean unit ball.  For \(1\leq k\leq n\), the \(k\)-th successive minimum
\(\lambda_k(\mathcal L)\) is the smallest radius \(r\) such that
\(r\mathcal B_2^m\) contains \(k\) linearly independent vectors of
\(\mathcal L\).  Equivalently,
\[
 \lambda_k(\mathcal L)
 :=\min\left\{r\geq0:
 \dim_{\mathbb R}\Span_{\mathbb R}
 \bigl(\mathcal L\cap r\mathcal B_2^m\bigr)\geq k
 \right\}.
\]
\end{definition}

\begin{definition}[Exact SIVP]
\label{def:exact-sivp}
Given a rational lattice basis \(\mathbf B\in\mathbb Q^{m\times n}\) and a
nonnegative rational squared threshold \(\Delta\), exact decision-SIVP asks
whether
\[
 \lambda_n\bigl(\mathcal L(\mathbf B)\bigr)^2\leq\Delta.
\]
\end{definition}
\fi

\Needspace{10\baselineskip}
We use the following formulation of exact decision-CVPP, following
Micciancio~\cite{Micciancio2001}.

\begin{definition}[Exact decision-CVPP]
\label{def:cvpp}
An exact decision-CVPP scheme consists of a
\emph{preprocessing function} $\mathsf{Pre}$, a
\emph{uniform polynomial-time decoder} $\mathsf{Dec}$, and a
polynomial $p$ with the following properties.
\begin{itemize}
 \item Given a rational lattice basis $\mathbf B$, the
       \emph{offline stage} outputs
       a string
       $\pi_{\mathbf B}:=\mathsf{Pre}(\mathbf B)$ satisfying
       $|\pi_{\mathbf B}|\leq p(|\langle\mathbf B\rangle|)$, where
       $\langle\mathbf B\rangle$ is the standard binary encoding of
       $\mathbf B$.  No running-time bound is imposed on
       $\mathsf{Pre}$.
 \item Given $\pi_{\mathbf B}$, a rational target $\mathbf t$, and a
       nonnegative rational squared threshold $\Delta$, the
       \emph{online stage}
       satisfies
       \[
        \mathsf{Dec}(\pi_{\mathbf B},\mathbf t,\Delta)=1
        \quad\Longleftrightarrow\quad
        \dist(\mathbf t,\mathcal L(\mathbf B))^2\leq\Delta.
       \]
       Its running time is polynomial in the total length of these inputs.
\end{itemize}
We say that exact decision-CVPP is solvable in polynomial time
on a lattice family with one specified basis for each lattice if a single such
scheme works for all these bases.
\end{definition}

\subsection{Algebraic Number Theory and Ideal Lattices}
\label{sec:simple-ideal-lattices}

Let $K$ be a number field of degree $n$, and let
$R:=\mathcal{O}_K$ be its ring of integers.

\paragraph{Coefficient embedding.}
Fix an algebraic integer $\theta\in R$ such that
$K=\mathbb{Q}(\theta)$, and let $F(z)\in\mathbb{Z}[z]$ be the monic
minimal polynomial of $\theta$.
Then $K\cong\mathbb{Q}[z]/(F(z))$.
Every element $a\in K$ has a unique representation
\[
  a=a_0+a_1\theta+\cdots+a_{n-1}\theta^{n-1},
  \qquad a_0,\ldots,a_{n-1}\in\mathbb{Q}.
\]
Its \emph{coefficient embedding} with respect to the
\emph{power basis}
$1,\theta,\ldots,\theta^{n-1}$ is
\[
  \coeff(a)
  :=(a_0,a_1,\ldots,a_{n-1})^\top\in\mathbb{Q}^n.
\]
If $I\subseteq R$ is a nonzero ideal, then
\[
  \mathcal{L}(I)
  :=\{\coeff(a):a\in I\}
\]
is a full-rank lattice in $\mathbb{R}^n$. {For $T\in K$, define}
\[
  {\dist(T,I)\mathrel{:=}}
  \operatorname{dist}\bigl(
    \coeff(T),\mathcal{L}(I)
  \bigr)
  \mathrel{{=}}
  \min_{a\in I}
  \|\coeff(T-a)\|_2.
\]

\paragraph{Canonical embedding.}
Let $\sigma_1,\ldots,\sigma_n:K\longrightarrow\mathbb{C}$ be all embeddings
of $K$ into $\mathbb{C}$. The \emph{canonical embedding} is defined by
\[
  \sigma(a)
  :=(\sigma_1(a),\ldots,\sigma_n(a))\in\mathbb{C}^n.
\]
Because the complex embeddings occur in conjugate pairs, $\sigma(K)$
lies in a real subspace of $\mathbb{C}^n$ of dimension $n$. If
$I\subseteq R$ is a nonzero ideal, then
\[
  \sigma(I):=\{\sigma(a):a\in I\}
\]
is a full-rank lattice in this real subspace. For $T\in K$, its distance
from $I$ under the canonical embedding is
\[
  \operatorname{dist}\bigl(\sigma(T),\sigma(I)\bigr)
  :=
  \min_{a\in I}\|\sigma(T-a)\|_2.
\]

In general, the coefficient and canonical embeddings induce different
Euclidean lengths. We now specialize to power-of-two
cyclotomic fields.

\paragraph{The power-of-two cyclotomic case.}
\label{sec:simple-cyclotomic}

Let $d$ be a power of two and put
$\zeta\mathrel{:=}e^{\pi i/d}$, a primitive
$2d$-th root of unity. We specialize to
\[
  K_d:=\mathbb{Q}(\zeta),
  \qquad
  R_d:=\mathcal{O}_{K_d}
       =\mathbb{Z}[\zeta]
       \cong\mathbb{Z}[y]/(y^d+1).
\]
In a quotient by a monic polynomial of degree \(r\), we identify
each element with its unique representative of degree less than \(r\).
Every $f\in K_d$ has a unique representation
\[
  f=f_0+f_1y+\cdots+f_{d-1}y^{d-1},
  \qquad f_0,\ldots,f_{d-1}\in\mathbb{Q},
\]
with coefficient vector
\[
  \coeff(f)
  =(f_0,\ldots,f_{d-1})^\top.
\]
The $d$ canonical embeddings of $K_d$ send $y$ to $\zeta^{2k+1}$ for
$k=0,\ldots,d-1$.
Accordingly, we write
\[
  \sigma_d(f)
  \mathrel{:=}\bigl(f(\zeta),f(\zeta^3),\ldots,f(\zeta^{2d-1})\bigr).
\]
For \(1\leq p\leq\infty\), we use
\(\|f\|_p:=\|\coeff(f)\|_p\) for the
\emph{coefficient \(\ell_p\)-norm} of
\(f\in K_d\).

We will use \emph{Young's convolution inequality} in the following form.
For every \(r\geq1\) and \(f,g\in\R[z]/(z^r+1)\), with coefficient norms
taken on their degree-\(<r\) representatives,
\[
 \|fg\|_2\leq\|f\|_1\|g\|_2.
\]
Indeed, multiplication by a monomial is a signed coordinate permutation, so
the inequality follows from the triangle inequality.

For \(g\in R_d\), write
\((g):=gR_d=\{gq:q\in R_d\}\) for the principal ideal generated by \(g\).
We write \(I\lhd R\) to mean that \(I\) is an ideal of \(R\).
For $0\neq g\in R_d$, let
\[
  \mathbf B_g
  :=\bigl[\coeff(g)\ \coeff(yg)\ \cdots\ \coeff(y^{d-1}g)\bigr].
\]
Since $R_d$ is a domain, these columns form a $\mathbb Z$-basis of
$\mathcal L((g))$. We refer to $\mathbf B_g$ as its
\emph{multiplication basis}.

We now formulate exact CVP on principal cyclotomic ideal lattices.
\begin{definition}[Principal cyclotomic CVP]
\label{def:principal-cyclotomic-cvp}
Given a power of two $d$, a nonzero generator $g\in R_d$, a target
$T\in K_d$, and a nonnegative rational squared threshold $\Delta$,
principal cyclotomic exact decision-CVP asks whether
\[
 {\dist(T,(g))^2\leq\Delta.}
\]
Given $d$, $g$, and $T$, principal cyclotomic exact search-CVP asks for
$\coeff(v)$, where $v\in(g)$ satisfies
\[
 {\|T-v\|_2=\dist(T,(g)).}
\]
\end{definition}

The inputs $g$ and $T$ are represented by $\coeff(g)$ and $\coeff(T)$,
respectively.  Under this representation, the corresponding standard
decision-CVP instance is $(\mathbf B_g,\coeff(T),\Delta)$.

In $K_d$, the canonical and coefficient
\(\ell_2\)-norms differ only by a fixed factor \(\sqrt d\).

\begin{lemma}
\label{lem:canonical}
For every $f\in K_d$,
\[\|\sigma_d(f)\|_2^2=d\|\coeff(f)\|_2^2.\]
\end{lemma}

\begin{proof}
For $0\leq r,s<d$, Fourier orthogonality gives
\[
  \sum_{k=0}^{d-1}\zeta^{(2k+1)(r-s)}
  =
  \begin{cases}
    d,&r=s,\\
    0,&r\neq s.
  \end{cases}
\]
Expanding the squared canonical norm therefore yields
\[\|\sigma_d(f)\|_2^2=d\sum_{r=0}^{d-1}|f_r|^2
=d\|\coeff(f)\|_2^2.\]
\end{proof}

\begin{remark}
Lemma~\ref{lem:canonical} implies that, for every nonzero ideal
$I\subseteq R_d$ and every target $T\in K_d$,
\[\operatorname{dist}\bigl(\sigma_d(T),\sigma_d(I)\bigr)^2
={d\dist(T,I)^2}.\]
Thus exact SVP and CVP under the coefficient and canonical embeddings are
equivalent after multiplying each squared threshold by $d$.
\end{remark}

\ifdefined\withModuleSIVP
\paragraph{Rank-two cyclotomic module lattices.}
For \(\mathbf v=(v_1,v_2)\in K_d^2\), define
\[
 \coeff_2(\mathbf v)
 :=\bigl(\coeff(v_1)^\top,\coeff(v_2)^\top\bigr)^\top
 \in\mathbb Q^{2d},
 \qquad
 \|\mathbf v\|_2:=\|\coeff_2(\mathbf v)\|_2.
\]
A free rank-two \(R_d\)-module in \(K_d^2\) has the form
\[
 \mathcal M=R_d\mathbf b_1+R_d\mathbf b_2,
\]
where \(\mathbf b_1,\mathbf b_2\in K_d^2\) are linearly independent over
\(K_d\).  Its \emph{coefficient lattice}
\[
 \mathcal L(\mathcal M)
 :=\{\coeff_2(\mathbf v):\mathbf v\in\mathcal M\}
\]
is a full-rank lattice in \(\mathbb R^{2d}\), called a
\emph{rank-two cyclotomic module lattice}.

The module analogue of the second successive minimum is defined as
\[
 \lambda_2^{K_d}(\mathcal M)
 :=\min\left\{r\geq0:
 \dim_{K_d}\Span_{K_d}
 \{\mathbf v\in\mathcal M:\|\mathbf v\|_2\leq r\}\geq2
 \right\}.
\]
For every nonzero \(\mathbf v\in\mathcal M\), the \(d\) vectors
\(y^j\mathbf v\) \((0\leq j<d)\) have equal norm and are
\(\mathbb R\)-linearly independent.  Hence
\[
 \lambda_2^{K_d}(\mathcal M)
 =\lambda_{d+1}\bigl(\mathcal L(\mathcal M)\bigr)
 =\lambda_{2d}\bigl(\mathcal L(\mathcal M)\bigr).
\]
For \(d=1\), this is the ordinary second successive minimum \(\lambda_2\).

\begin{definition}[Cyclotomic module-SIVP]
\label{def:module-sivp}
Given a power of two \(d\), an \(R_d\)-basis of a free rank-two module
\(\mathcal M\subseteq K_d^2\), and a nonnegative integer squared threshold
\(\Gamma\), exact decision-module-SIVP asks whether
\[
 \bigl(\lambda_2^{K_d}(\mathcal M)\bigr)^2\leq\Gamma.
\]
\end{definition}
{The input basis vectors are represented by their coefficient
vectors under \(\coeff_2\).}
For \(\mathbf v=(v_1,v_2)\in K_d^2\), applying
Lemma~\ref{lem:canonical} componentwise gives
\[
 \|\bigl(\sigma_d(v_1),\sigma_d(v_2)\bigr)\|_2^2
 =d\|\mathbf v\|_2^2.
\]
Thus the coefficient norm and the componentwise canonical norm differ by the
fixed factor \(\sqrt d\).
\fi

\subsection{Exact Cover by 3-Sets (X3C).}
Fix a positive multiple \(m\) of \(3\), and let
\(\mathcal C\mathrel{:=}\{C_1,\ldots,C_n\}\) be a collection of
three-element subsets of \([m]\) (called \emph{triples}).
The \emph{incidence matrix}
$\mathbf{H} \in \{0,1\}^{m \times n}$ of the instance is defined by
\[
H_{ij} \mathrel{:=}
\begin{cases}
1 & \text{if }i\in C_j,\\[2pt]
0 & \text{otherwise.}
\end{cases}
\]
An \emph{exact cover} is a vector $\boldsymbol{\xi} \in \{0,1\}^n$ satisfying
the binary linear system
\[
\mathbf{H}\boldsymbol{\xi} = \mathbf{1}_m,
\]
where $\mathbf{1}_m \in \{0,1\}^m$ denotes the all-ones vector.
X3C asks whether such an exact cover exists.
Equivalently, the selected subfamily
$\mathcal{C}' \mathrel{:=} \{C_j : \xi_j = 1\} \subseteq \mathcal{C}$ is pairwise
disjoint and covers \([m]\) exactly once:
\[
\bigcup_{C_j \in \mathcal{C}'} C_j = [m],
\qquad
C_j \cap C_{j'} = \varnothing
\;\text{ for all distinct } C_j, C_{j'} \in \mathcal{C}'.
\]

It is well known that X3C is \(\NP\)-complete~\cite{GareyJohnson}.
